\section{Introduction}
\label{sec:intro}

Cities concentrate the risks of a warming climate, and the form of the city is
increasingly understood not as a passive backdrop to those risks but as one of
their controls \citep{hamin2009UrbanForm, reckien2018CitiesPlanning}. Recent work shows that urban morphology governs land-surface
temperature through both marginal and interaction effects, so that the thermal
behaviour of a neighbourhood cannot be read off its density alone
\citep{wangZ2025MarginalInteraction}; that heat exposure and the social capacity
to cope with it are spatially uneven and frequently decoupled across urban forms
\citep{turner2025Guhvi, iqbal2025HeatStress}; and that heat and pluvial flooding
act as compound, interacting hazards rather than independent ones
\citep{imroz2025IntegratedAssess}. Distributive justice and holistic health
adaptation require urban climate planning that explicitly targets these form-based
realities \citep{chu2025EvaluatiCities, o2026CitiesNeed}. If form shapes risk,
then adaptation that ignores form will misallocate effort.

In parallel, the measurement of urban form has matured rapidly. Urban
morphometrics has moved from case-study description to a reproducible,
unsupervised science capable of classifying built fabric at national and even
global scale, using consistent spatial units and open libraries
\citep{fleischmann2020Morphologic, fleischmann2022MethodolFoundati,
araldi2025MultiLevel, debray2025UniversalPatterns}, with interpretable, citizen-guided,
and deep-learning variants proliferating \citep{vartholomaios2025Detection,
fang2024Urbanclassifier, wangJ2024FigureGround, campomanes2026WhoseCity}. The
street-network side is similarly consolidated around reproducible open tooling
\citep{boeing2017Osmnx, boeing2025ModelingNetworks}. This computational ecosystem includes Python-based libraries like \texttt{momepy} and \texttt{tessPy} \citep{tesspy_zenodo} alongside established GIS-integrated software such as the Spatial Design Network Analysis (\texttt{sDNA}) engine and the Urban Multi-scale Environmental Predictor (\texttt{UMEP}) plugin suite for QGIS \citep{cooper2020sDNA, lindberg2018UMEP}. The field can now say,
rigorously and at scale, what kind of fabric a place is.

Yet these two trajectories rarely meet at the scale where adaptation is actually
designed, financed and lived: the street and the walkable neighbourhood.
Climate-resilience assessment remains dominated by administrative units, coarse
grids and Local Climate Zone tiles. Heat-vulnerability indices, even excellent
open, transferable and multi-hazard ones, are typically resolved at the neighbourhood
or census scale and are not mechanistically tied to morphology
\citep{turner2025Guhvi, rupp2026WhereHeat}. Multiscale spatial modelling in
Istanbul further shows that the diagnosis of heat vulnerability changes with
analytical scale and spatial specification \citep{eminoglu2026Diagnosing}.
Morphometric typologies, conversely, remain largely descriptive and are seldom
connected to hazard outcomes or to adaptation decisions
\citep{araldi2025MultiLevel, debray2025UniversalPatterns}.
The consequence is a practical blind spot: two areas of similar density and land
use can be assigned the same risk class while producing materially different heat,
accessibility and vulnerability outcomes, a contrast documented empirically but
obscured by coarse-unit assessment \citep{wangZ2025MarginalInteraction,
iqbal2025HeatStress}.

This study addresses that blind spot with an open-source, auditable, QGIS-based
typomorphological workflow for the \.{I}zmir functional urban region, a coastal
metropolitan region with demonstrated and recent exposure to extreme heat,
pluvial flooding and coastal hazard \citep{canguzel2024ClimateChange}. The
workflow couples street-scale morphometrics with multi-hazard exposure and social
vulnerability to produce transparent, fabric-specific adaptation priorities. Its
pluvial-susceptibility output is a screening overlay rather than one of the five
priority axes; hydrodynamic validation is required before it can enter the
multi-objective ranking. Its
methodological core is not the individual tools but a catchment-radius
cross-attribution rule that renders movement potential, built-form intensity and
hazard exposure commensurable within a single analytical unit, so that the three
are integrated rather than reported as parallel, unlinked layers. The workflow is
deliberately distinguished from form-only typologies, which stop at description,
and from administrative resilience indices, which miss the street scale at which
form is produced and experienced.

The study distinguishes \emph{urban fabric} in three registers: type,
measurement and mechanism, turning an intuition into a testable proposition. Its
integration unit is the 250\,m grid cell, with tessellation-cell morphometrics and
network metrics cross-attributed through 400/800\,m network service-area reaches
to make morphometric, network and hazard information commensurable. Within that
unit, an explainable gradient-boosting model with SHAP attributes measured
land-surface temperature to morphological mechanisms
\citep{wangZ2025PlainPlateau, liu2026Nonlinear, li2026BlockGrid}, while Pareto
optimization with TOPSIS and entropy/Monte-Carlo weight robustness ranks
trade-off-aware adaptation priorities across five need axes
\citep{zhang2024MultiObjective, zhu2025SpongeCities}. These stages are
complementary rather than nested: SHAP diagnoses why a fabric is hot and which
intervention fits it, while the optimizer ranks where to act using measured
land-surface temperature rather than a SHAP-derived heat axis. The resulting
open, parameter-logged implementation covers the full census of 3{,}777 urban
cells and runs Moran/LISA, Getis-Ord $G_i^{*}$ and a Gini equity measure in the
PlanX QGIS suite.

Four questions structure the analysis. The first asks whether the seven
provisional fabric types are distinguishable on measured morphometric,
accessibility and hazard-exposure profiles. The second asks whether the
resulting classifications are stable across grid resolutions, tested by
re-deriving the clustering at 500\,m and comparing it to the 250\,m baseline
via the Adjusted Rand Index. The third asks which morphological mechanism,
among those captured by the 29-indicator feature vector, carries the greatest
attributable effect on measured summer land-surface temperature within each
fabric type, as identified by SHAP decomposition of a spatial-block-validated
gradient-boosting model. The fourth asks where adaptation priorities fall when
the five need axes are treated simultaneously as a multi-objective problem:
which fabrics are Pareto-dominated (unambiguous intervention candidates) versus
on the trade-off frontier, and how robust is the TOPSIS ranking to entropy
weighting and Monte-Carlo weight perturbation?
